\documentclass[10pt, a4paper, landscape]{article}

% --- Paquetes Esenciales ---
\usepackage[utf8]{inputenc}
\usepackage[spanish]{babel}
\usepackage{listings}
\usepackage{xcolor}
\usepackage{multicol}
\usepackage{geometry}
\usepackage{fancyhdr}
\usepackage{tocloft}

% --- Configuración de Página (Márgenes tipo ICPC) ---
\geometry{top=1.2cm, bottom=1.2cm, left=1cm, right=1cm}
\columnseprule = 0.4pt % Línea divisoria entre columnas

% --- Configuración de Colores (Estilo Morass/KACTL) ---
\definecolor{codegray}{rgb}{0.5,0.5,0.5}
\definecolor{codeblue}{rgb}{0.0,0.0,0.8}
\definecolor{codegreen}{rgb}{0,0.6,0}

% --- Estilo de Código (Listings) ---
\lstset{
    language=C++,
    basicstyle=\ttfamily\scriptsize, % Fuente pequeña para que quepa más
    keywordstyle=\color{codeblue}\bfseries,
    commentstyle=\color{codegray}\itshape,
    stringstyle=\color{codegreen},
    breaklines=true,                % Salto de línea automático
    showstringspaces=false,
    tabsize=2,
    frame=single,                   % Recuadro alrededor del código
    rulecolor=\color{black!30},
    numbers=left,                   % Números de línea a la izquierda
    numberstyle=\tiny\color{codegray},
    stepnumber=1,
    numbersep=5pt,
    firstnumber=1
}

% --- Encabezado y Pie de Página ---
\pagestyle{fancy}
\fancyhf{}
\fancyhead[L]{Pumas Sangrientos - Universidad Nacional Autónoma de México}
\fancyhead[R]{Página \thepage}

\begin{document}

% --- Título Compacto ---
\begin{center}
    {\Huge \textbf{Pumas Sangrientos: Team Reference Document}} \\
    \textit{ICPC Training - 2024}
\end{center}

\begin{multicols*}{2} % El '*' mantiene las columnas llenas antes de pasar a la siguiente
    \tableofcontents
    \vspace{1cm}

    % --- SECCIÓN DE EJEMPLO ---
    % Nota: El script de Python debería insertar las secciones aquí
    \section{Data Structures}
    
    \subsection{Segment Tree (Lazy Propagation)}
    % Aquí se enlaza el archivo local
    % \lstinputlisting{code/data_structures/seg_tree.cpp}
    \begin{lstlisting}
    // Ejemplo rápido de cómo se verá el código
    void build(int node, int start, int end) {
        if(start == end) {
            tree[node] = A[start];
            return;
        }
        int mid = (start + end) / 2;
        build(2*node, start, mid);
        build(2*node+1, mid+1, end);
        tree[node] = tree[2*node] + tree[2*node+1];
    }
    \end{lstlisting}

    \section{Graphs}
    \subsection{Dinic's Algorithm}
    % \lstinputlisting{code/graphs/dinic.cpp}

\end{multicols*}
\end{document}